% ============================================================
% SECCION 10 - CONCLUSIONES (Roselyn Sanchez - PE5)
% Cinco conclusiones integradoras, una por unidad del curso,
% ancladas en el PFC AgroMoreira y no aplicables a cualquier
% proyecto.
% ============================================================
\section{Conclusiones}
\label{sec:conclusiones}

Se presentan cinco conclusiones integradoras, una por cada unidad del curso, que vinculan lo aprendido con el resultado final del SGA de AgroMoreira.

\subsection*{Unidad I --- Fundamentos y elicitation de requisitos}

La combinaci\'on de entrevista semiestructurada, encuesta digital y observaci\'on directa result\'o decisiva para un dominio donde los datos no existen todav\'ia: la finca registraba en bit\'acoras f\'isicas y comunicaba verbalmente sus labores, y ninguna t\'ecnica por s\'ola habr\'ia revelado tanto la brecha entre perfiles (un administrador que decide con datos y jornaleros que ejecutan con instrucciones) como la amenaza fitosanitaria real ---moniliasis, sigatoka y una enfermedad sin identificar--- que luego justific\'o el componente IA-02. La lecci\'on metodol\'ogica es que elicitar es un proceso sistem\'atico con t\'ecnicas seleccionadas por acceso al stakeholder y por el tipo de informaci\'on buscada \cite{pohl2025}, no una conversaci\'on informal; el marco conceptual de SWEBOK v4.0 permiti\'o ordenar evidencia (EV-01 a EV-07) con trazabilidad hacia atr\'as verificable desde el primer d\'ia \cite{swebok2024}.

\subsection*{Unidad II --- Modelado y especificaci\'on}

Pasar de 15 requisitos brutos con lenguaje aspiracional (``deber\'ia permitir'') a un cat\'alogo de 57 requisitos especificados con plantilla de ocho atributos solo fue posible cuando el equipo acept\'o que un requisito es una afirmaci\'on verificable sobre el sistema y no una intenci\'on. La conformancia con ISO/IEC/IEEE 29148:2018 dej\'o de ser un requisito de forma cuando se tradujo en decisiones concretas: vistas derivadas referenciando el requisito base, criterios Dado--Cuando--Entonces y umbrales num\'ericos en cada RNF \cite{iso29148}. El modelado UML (contexto, casos de uso, clases refinadas, secuencias, actividad y estados) cumpli\'o su funci\'on real: detectar temprano que RF-28 a RF-31 exig\'ian casos de uso nuevos, evitando requisitos hu\'erfanos antes de llegar a la matriz final. Las listas de verificaci\'on del temario IREB CPRE sirvieron como instrumento operativo de revisi\'on, no como teor\'ia de repaso \cite{ireb2026}, y el encadre de procesos de ISO/IEC/IEEE 15288 e ISO/IEC/IEEE 12207 orden\'o la definici\'on de requisitos dentro de un ciclo de vida coherente \cite{iso15288,iso12207}.

\subsection*{Unidad III --- Validaci\'on y calidad del producto documental}

La inspecci\'on de la PE4 y la re-inspecci\'on de la PE5 demostraron emp\'ricamente que la calidad de una ERS se mide, se corrige y se vuelve a medir: de los defectos residuales iniciales (0,119 por requisito) se lleg\'o a 0,024 tras cerrar DR-01 a DR-05, bajo el umbral de 0,05 exigido por M6. El hallazgo m\'as valioso no fue ning\'un defecto individual sino el patr\'on detr\'as de ellos: casi todos proven\'ian de redactar r\'apido sin umbral o de duplicar contenido entre secciones, exactamente las dos pr\'acticas que la retrospectiva manda detener. Verificar que las correcciones no introdujeran nuevos defectos exigi\'o conteos reproducibles por script y comparaci\'on por hash entre corridas, convirtiendo la validaci\'on documental en un experimento controlado con aritm\'etica visible \cite{wohlin2012,iso25010}.

\subsection*{Unidad IV --- Gesti\'on, trazabilidad y trabajo en equipo}

La cadena end-to-end Fuente $\rightarrow$ RF $\rightarrow$ CU $\rightarrow$ Clase $\rightarrow$ Estado $\rightarrow$ BDD $\rightarrow$ CP dej\'o de ser un esquema te\'orico cuando la matriz oblig\'o a decidir qu\'e hacer con cada fila rota: enlazar, eliminar o justificar por escrito, sin excepciones. La experiencia confirma el problema cl\'asico de la trazabilidad reportado desde Gotel y Finkelstein: mantenerla es costoso si se construye despu\'es, y barata si se registra en el momento \cite{gotel1994}; la literatura reciente sobre trazabilidad pre-especificaci\'on ayud\'o a distinguir qu\'e eslabones corresponden a la evidencia y cu\'ales a los artefactos derivados \cite{mucha2024}. La sincronizaci\'on backlog--ERS (cada RF con historia asociada y cada historia con RF de origen) cerr\'o la brecha entre la gesti\'on \'agil del tablero y el documento formal, resolviendo con disciplina el desaf\'io conocido de combinar ambos mundos \cite{hoyxu2023}. En el plano del equipo, el reparto por pasos con defensa cruzada ---cada integrante puede explicar el bloque de otro--- fue la decisi\'on de gesti\'on que sostuvo la calidad individual medida por el FAI.

\subsection*{Unidad V --- Inteligencia artificial, m\'etricas y cierre}

Especificar los dos componentes de IA (IA-01 recomendador de manejo agr\'icola; IA-02 detector de plagas por imagen) oblig\'o a aplicar un paradigma distinto al del resto del cat\'alogo: el comportamiento se induce de datos y se eval\'ua con m\'etricas probabil\'isticas, por lo que cada requisito naci\'o con umbral, unidad y m\'etodo de comprobaci\'on ---exactitud/F1 m\'inima sobre conjunto de prueba, latencia percentilada, brechas m\'aximas de equidad entre cultivos, parcelas y dispositivos--- siguiendo las recomendaciones de ingenier\'ia de requisitos para aprendizaje autom\'atico \cite{vogelsang2019} y la evidencia de c\'omo los equipos reales luchan por operacionalizar estos RNF \cite{habibullah2023,ahmad2023}. Los bloques de datos, equidad, explicabilidad y monitoreo por deriva transformaron obligaciones legales en requisitos trazables: la clasificaci\'on de riesgo del Reglamento (UE) 2024/1689 defini\'o supervisi\'on humana y transparencia \cite{eu20241689}, y la LOPDP ecuatoriana con su reglamento fij\'o base legal, minimizaci\'on y plazos de conservaci\'on para los registros productivos tratados \cite{lopdp2021,decreto904}. Finalmente, la auditor\'ia con las seis m\'etricas de ISO/IEC 25010:2023 cerr\'o el c\'irculo de la asignatura: todo valor reportado tiene conteo exhibible, toda correcci\'on tiene par antes/despu\'es y toda afirmaci\'on de esta defensa tiene un artefacto concreto que la respalda \cite{iso25010}.
